% ECONOMICS_PROBLEM: {"id": "D72-631", "title": "Causal effects of money in politics", "jel_code": "D72", "source_id": "OP-0122"}
% !TeX program = xelatex
\documentclass[11pt,a4paper]{article}
\usepackage[margin=23mm]{geometry}
\usepackage{fontspec}
\setmainfont{Latin Modern Roman}
\usepackage{amsmath,amssymb,mathtools}
\usepackage{xcolor,enumitem,booktabs,longtable,array,xurl,hyperref}
\definecolor{muted}{HTML}{53636C}
\hypersetup{colorlinks=true,urlcolor=blue,linkcolor=blue}
\setlength{\parindent}{0pt}
\setlength{\parskip}{5pt}
\newcommand{\R}{\mathbb R}
\newcommand{\E}{\mathbb E}
\newcommand{\N}{\mathbb N}
\newcommand{\ind}{\mathbf 1}
\newcommand{\Var}{\operatorname{Var}}
\newcommand{\Cov}{\operatorname{Cov}}
\newcommand{\supp}{\operatorname{supp}}
\DeclareMathOperator*{\argmin}{arg\,min}
\DeclareMathOperator*{\argmax}{arg\,max}
\newcommand{\entrymeta}[3]{{\sffamily\small\color{muted}\textbf{\texttt{#1}}\quad|\quad #2\par #3\par}}
\newcommand{\Problemref}[1]{\texttt{#1}}
\newcommand{\JELref}[2]{\texttt{#2} (JEL \texttt{#1})}
\begin{document}
\section*{Causal effects of money in politics}
\textbf{Problem ID:} \texttt{D72-631}\quad \textbf{JEL:} \texttt{D72}\par
% BEGIN ECONOMICS STATEMENT
\label{problem:OP-0122}
\label{atlas:prob:Q2}
{\sffamily\small\color{muted}\textbf{Q2} | Political economy | Editorial research problem family\par}
\subsection*{Definitions and assumptions}
Fix elections and a subsequent policy horizon $h$. Let $a=(a^{\rm campaign},a^{\rm lobby})\in\mathcal A$ be a feasible vector of spending interventions specifying the donor, recipient, issue, message, timing and disclosure regime. For jurisdiction $j$, let $V_j(a)$ denote election outcomes and $Y_j(a)$ a specified vector of legislative, procurement or regulatory decisions after the election. These potential outcomes incorporate opponents' responses under a fixed equilibrium selection rule. Let $\mathcal M$ specify donor preferences, candidate alignment, strategic spending, policy determination and the data-observation process. Outcomes are integrable; policy components use declared units. A randomized spending encouragement $Z_j$ is a distinct variable from realized spending.
\subsection*{Formal research question}
Determine the identified set for
\[
\theta_Y(a,a_0)=\mathbb E[Y_j(a)-Y_j(a_0)],\qquad
\theta_V(a,a_0)=\mathbb E[V_j(a)-V_j(a_0)],
\]
under an explicit assignment or quasi-experimental design and law class. When an instrument $Z_j$ is used, state the exclusion, independence, relevance and compliance assumptions, and identify the resulting local target rather than automatically equating it to $\theta_Y$. If the question is about influence holding the election winner fixed, define that intervention separately and determine whether the data and model identify it.
\subsection*{Interpretation and scope}
An effect on advertising exposure, votes or election victory is different from an effect on the winner's policy. Observed donations may reflect shared ideology, expected electoral success or anticipated access. A mediation decomposition into candidate selection and post-election persuasion requires additional assumptions, including a defensible intervention on the mediator; ordinary regressions controlling for the winner do not provide it. The research target includes agenda setting and administrative decisions when those are part of the prespecified policy vector.
\subsection*{Source audit and status}
All three atlas author-year citations are verified. [S1] and [S2] provide foundations or reviews; [S3] identifies advertising effects on elections and is only an indirect source for policy buying. The exact target is editorial. Existing causal estimates in specific settings are acknowledged; a uniform law linking money to policy has not been established by these sources. This is an empirical identification agenda rather than an unresolved theorem.
\subsection*{References}
{\footnotesize\raggedright
\begin{enumerate}[label={\textbf{[S\arabic*]}},leftmargin=3em]
\item Stephen Ansolabehere, John M. de Figueiredo, and James M. Snyder Jr. (2003). \emph{Why is There so Little Money in U.S. Politics?}. Journal of Economic Perspectives 17(1), 105--130.\par
\textit{Locator and source role:} Abstract and review of political contributions; spending levels and candidate selection do not by themselves identify policy purchases.\par
\url{https://www.aeaweb.org/articles?id=10.1257/089533003321164976}
\item John M. de Figueiredo and Brian Kelleher Richter (2014). \emph{Advancing the Empirical Research on Lobbying}. Annual Review of Political Science 17, 163--185.\par
\textit{Locator and source role:} Abstract and research-agenda discussion; bibliographic match verified. NBER Working Paper 19698 is the 2013 version.\par
\url{https://www.annualreviews.org/doi/10.1146/annurev-polisci-100711-135308}
\item Jörg L. Spenkuch and David Toniatti (2018). \emph{Political Advertising and Election Results}. Quarterly Journal of Economics 133(4), 1981--2036.\par
\textit{Locator and source role:} Abstract and advertising exposure design; election effects rather than a causal effect of contributions on subsequent policy.\par
\url{https://academic.oup.com/qje/article-abstract/133/4/1981/4993157}
\end{enumerate}
}

\subsection*{Common conventions: Source mathematical conventions}

Unless an entry states otherwise, agent and firm sets are finite. State and action spaces are measurable, strategies and outcome maps are measurable, probability laws are specified, and expectations exist for the targets under discussion. Infinite-horizon discount factors lie in $(0,1)$ and discounted objectives are finite. Norms, tolerances and complexity input sizes must be specified for computational comparisons. Constants or primitives that appear locally are inputs, not empirical facts asserted by this catalogue.

\textbf{Identification.} A statistical model $m\in\mathcal M$ implies an observable law $P_m$ and target $\theta(m)$. At law $P$, its identified set is
\[
\Theta(P;\mathcal M)=\{\theta(m):m\in\mathcal M,\ P_m=P\}.
\]
Point identification holds when this set is a singleton, or equivalently when $P_{m_1}=P_{m_2}$ implies $\theta(m_1)=\theta(m_2)$ for all compatible models. Sharp bounds mean bounds attained, or approached when the boundary is not attained, by compatible models. Writing this set is a definition, not a solution: the substantive task is to establish informative restrictions and derive or estimate the set. An unrestricted-confounding model may give only trivial bounds.

\textbf{Inference.} Identification, estimability and valid uncertainty quantification are separate questions. Fix $\alpha\in(0,1)$. A confidence procedure $C_n$ over a declared law class $\mathcal P$ has asymptotic uniform coverage at level $1-\alpha$ when
\[
\liminf_{n\to\infty}\inf_{P\in\mathcal P}\Pr_P\{\theta(P)\in C_n\}\geq1-\alpha.
\]
For set-identified targets the coverage event must be defined separately, for example coverage of the full identified set. If an entry uses identification-indexed models instead of uniquely indexed laws, the same coverage convention is applied over those models. Pointwise limits do not establish uniform validity, and asymptotic validity does not guarantee finite-sample precision. Sample sizes, dependence structures and weak-identification sequences are part of each application.

\textbf{Counterfactuals.} $Y_i(a)$ is a potential outcome under a specified intervention, including its timing, implementation and versions. It is not $Y_i$ conditional on voluntarily choosing $a$. A full intervention vector permits interference; shorthand scalar treatment notation does not silently impose no interference. Assignment, selection, attrition, anticipation and measurement must be specified. A claim about an instrument's local effect is not automatically a population policy effect.

\textbf{Economic equilibrium.} A counterfactual model includes best responses, constraints, feasibility, market clearing, state transitions and expectations consistent with the stated information. When several equilibria exist, either a declared selector is an input or the target is set-valued. Comparisons across policy regimes must use the stated selection convention. Reduced-form relationships are not presumed invariant after an equilibrium-changing intervention.

\textbf{Optimization and welfare.} Optimization is over the stated feasible set. If attainment is not established, $\sup$ or $\inf$ and $\varepsilon$-optimal choices replace an asserted maximizer or minimizer. For example, with value $V=\sup_{a\in\mathcal A}W(a)<\infty$, an $\varepsilon$-optimal set is $\{a\in\mathcal A:W(a)\geq V-\varepsilon\}$ for $\varepsilon>0$. Benefits, costs, risk and health or environmental outcomes can be aggregated only using declared units and conversion weights. Interpersonal utility weights, the treatment of fiscal transfers and foreign profits, discounting, and the behavioral welfare criterion are explicit inputs. A comparative-static derivative additionally requires existence and appropriate differentiability of the selected counterfactual.

\textbf{Sources.} Local labels [S1], [S2], and so forth restart in every question. Locators identify the relevant abstract, model, section or result at the level actually checked; a bibliographic or abstract-level verification is not represented as inspection of an entire proof. Working papers are distinguished from later publications and changing versions. Source summaries are paraphrased. All URL checks and editorial formulations refer to the audit date stated above.
% END ECONOMICS STATEMENT
\end{document}
